\documentclass[conference]{IEEEtran}
\usepackage[utf8]{inputenc}
\usepackage[T1]{fontenc}
\usepackage{cite}
\usepackage{amsmath,amssymb}
\usepackage{graphicx}
\usepackage{booktabs}
\usepackage{url}

\title{Verifier-Guided Prompting for Intent-to-Configuration Translation: A Preregistered NetOps Study}

\author{
\IEEEauthorblockN{Autonomous AI Researcher}
\IEEEauthorblockA{CNSM 2026 Agentic AI Researcher Track}
}

\begin{document}

\maketitle

\begin{abstract}
We preregistered and executed a paired confirmatory trial to test whether constrained prompting plus a deterministic verifier-guided generate\ensuremath{\rightarrow}verify\ensuremath{\rightarrow}repair loop (one allowed repair) increases deterministic-verifier semantic-correctness when translating operator natural-language intents into low-level device configurations. Design: N=150 synthetic intents sampled from a deterministic generator, shared\_initial\_candidate generation semantics, model = gpt-5-mini, deterministic validator = deterministic\_netops\_validator\_v5, one initial generation per intent shared across baseline and guarded conditions, guarded pipeline permitted <=1 repair call. Results (artifact-grounded): baseline = 149/150 (99.333\%); guarded = 150/150 (100.0\%); paired absolute difference = 0.006667 (\ensuremath{\approx}0.67 percentage points); exact McNemar two-sided p = 1.0; 95\% bootstrap percentile CI = [0.00, 0.02] (10,000 resamples, seed = 1729). Provider accounting: completed episodes = 300; model\_calls\_used = 151; repair-stage invocations = 1. Interpretation: on this preregistered synthetic corpus and under shared\_initial\_candidate semantics the guarded workflow produced a numerically negligible and statistically non-significant improvement over a near-ceiling baseline. We emphasize the methodological lesson that preregistered NetOps trials require baseline-calibration pilots to avoid ceiling effects that invalidate preregistered power assumptions. All execution and analysis artifacts are archived in the supplied execution bundle referenced by the execution manifest.
\end{abstract}

\section{Introduction}
Natural-language intent interfaces aim to reduce operator burden by letting humans express desired network changes as free text that is automatically translated into executable device configurations. Prior work documents that single-pass LLM outputs often contain syntactic errors, omissions, or semantic violations that can yield unsafe or unavailable states if applied directly [1][2][3]. A commonly proposed defense is a generate\ensuremath{\rightarrow}verify\ensuremath{\rightarrow}repair architecture in which an LLM produces a candidate, a deterministic verifier checks intent-level invariants or formalized constraints, and an automated repair (or human gate) corrects violations before execution [4][5]. Security studies of tool-augmented LLM agents further motivate deterministic gating because unchecked textual claims can become harmful actions in multi-tool workflows [6].

Research question. After an autonomous literature review and preregistration we posed a confined confirmatory question: does constrained prompting (structured templates with explicit semantic slots) combined with deterministic verifier-guided iterative feedback materially increase verifier-level semantic correctness of NL\ensuremath{\rightarrow}configuration translation versus an unconstrained baseline under a paired design? The trial (study\_cand\_2026\_b2\_v1) was preregistered and frozen before any model outputs were observed.

Contributions. (1) A fully preregistered, artifactized paired confirmatory trial measuring deterministic-verifier pass/fail on a programmatically generated synthetic corpus (N=150 paired intents). (2) Artifact-backed outcomes showing that, under shared\_initial\_candidate semantics and a one-repair budget, the guarded pipeline raised verifier pass rate from 149/150 to 150/150 (+0.67 pp), a numerically negligible and statistically non-significant improvement (exact McNemar p = 1.0; 95\% bootstrap CI = [0.00, 0.02]). (3) A methodological recommendation that preregistered NetOps trials include baseline-calibration pilots and explicitly specify generation semantics (shared\_initial\_candidate vs independent sampling) because semantics materially change the estimand and interpretation.

\section{Related Work}
Related literature falls into three strands. First, intent-to-configuration work demonstrates that converting free-text intents into verifier-sound artifacts is difficult and reports modest semantic-correctness on benchmarks, motivating post-generation checks [1][2]. Second, generate\ensuremath{\rightarrow}verify\ensuremath{\rightarrow}repair architectures---ranging from verifier-in-the-loop fine-tuning to policy-guided verifier feedback---have shown correctness improvements for infrastructure-as-code and configuration repair, but reported gains vary with baseline model competence, repair budget, and sampling semantics [4][7]. Third, security and systems analyses have demonstrated claim-to-action vulnerabilities where multi-tool LLM pipelines can convert false textual claims into concrete harmful actions, arguing for auditable gating and verifier-aware planning [6][8].

How this study situates. Our trial operationalizes a narrowly specified guarded architecture and emphasizes preregistration, paired inference, and deterministic-verifier endpoints. Key contrasts with prior reports that observed larger guarded-vs-baseline gains are: (a) shared\_initial\_candidate execution semantics (one initial sample shared across conditions), (b) a strict one-repair budget per intent, (c) a single tested model (gpt-5-mini), and (d) a synthetic verifier-capable task bank restricted to small topologies. These design choices reduce the maximal observable marginal benefit of the guarded repair under our estimand and explain why previously reported larger gains can coexist with our bounded result [4][7].

\section{Methodology}
Overview and preregistration. We preregistered study\_cand\_2026\_b2\_v1 and froze an analysis plan prior to any model calls. The design is paired-binary: each sampled intent yields one shared initial candidate used to evaluate both the baseline and guarded conditions (shared\_initial\_candidate semantics). The primary estimand was the paired\_success\_rate\_difference\_guarded\_minus\_baseline and the confirmatory test was an exact two-sided McNemar on the paired contingency table. The preregistration specified N=150 confirmatory intents, inclusion/exclusion rules, a one-repair budget for guarded episodes, a deterministic retry policy for provider timeouts (one retry), and an analysis plan that included a 10,000-resample paired bootstrap (bootstrap\_seed = 1729) to produce percentile CIs. These preregistered elements are archived; the registered preregistration SHA-256 is recorded in the execution manifest.

Task generation, intent families, and prechecks. Confirmatory tasks were produced deterministically by deterministic\_netops\_task\_generator\_v5. The generator emits intent payloads with: (i) an explicit natural-language intent string, (ii) an initial\_state and expected\_final\_state, (iii) an ordered required\_sequence of field-level operations (e.g., shutdown\ensuremath{\rightarrow}mtu\ensuremath{\rightarrow}restore), (iv) policy constraints (allowed\_touched\_fields, protected\_interfaces, final\_group\_constraints), and (v) a declared difficulty label. The confirmatory corpus was balanced across four intent families: safe sequence reachability/isolation, ACL-like access changes, VLAN/MTU sequence manipulations, and uplink switchover patterns. Topologies were restricted to verifier-capable small networks (4--32 nodes) to ensure deterministic validator coverage. Prior to model calls we ran deterministic unit-tests on the mapping/transformation code and a rule-based ambiguity detector; items flagged as ambiguous or failing parser/unit-tests were excluded per the preregistered exclusion rules. In this execution none of the confirmatory items were excluded.

Execution adapter, sampling semantics, and explicit baseline interpretation. The hosted\_netops\_gvr\_v1 adapter enforced the shared\_initial\_candidate semantics required by the preregistration: for each intent the adapter performed exactly one initial model generation (the shared candidate) and used that identical candidate when scoring both the baseline and guarded conditions. The guarded pipeline could issue at most one additional repair call only if the deterministic validator failed the shared initial candidate. Because of this architecture, the baseline success rate measures the validity of the single initial sampled candidate; the guarded vs baseline contrast therefore quantifies the marginal effect of a single allowed repair applied to the identical sampled candidate rather than the effect of independently re-sampling first-pass candidates per condition. We include this precise sampling-to-estimand mapping in Methods to eliminate interpretive ambiguity about whether different first-pass generations were drawn for each condition.

Model configuration and nondeterminism. Declared model: gpt-5-mini (provider record: openai\_responses). The archived execution/model\_configuration.json records declared\_model\_version = "gpt-5-mini" and explicitly records temperature = null (unset). Importantly, the adapter and preregistration did not enforce a numeric temperature value or a fixed random seed; therefore model outputs may be nondeterministic across independent API sessions. We state this here in Methods and repeat it in the Disclosure Statement to comply with preregistration/runtime-parameter reporting requirements.

Repair budget, session isolation, and provider auditing. The guarded repair policy allowed at most one repair invocation per intent and the adapter maintained strict session isolation: each model call ran in a fresh API session (no cross-call state reuse). Provider auditing tracked hosted\_model\_call\_event\_count and cache statistics; the deterministic reconciliation artifact reports completed\_hosted\_model\_call\_event\_count = 151, cache\_miss\_count = 151, and stage\_counts = \{shared\_initial\_generation: 150, repair: 1\}. The adapter implementation enforces the preregistered retry rule (one retry after transient provider failure). All provider-call metadata and per-call runtime parameters were archived for audit.

Deterministic validator, scoring rule, and diagnostic trace structure. We used deterministic\_netops\_validator\_v5 (validator\_version = 5.0) as a deterministic oracle for the primary binary endpoint. The validator implements a full\_intent\_constraint\_validation rule that: (a) parses the LLM-produced candidate into normalized command sequences, (b) computes parsed\_command\_count and compares against required\_command\_count derived from the task's required\_sequence, (c) enforces workflow and group constraints (e.g., make-before-break, must\_be\_down\_for\_fields), and (d) returns a boolean valid flag plus structured diagnostics (violation\_codes, observed\_touched\_field\_count, normalized\_configuration, final\_state snapshot). Scoring for the confirmatory endpoint used the validator's boolean pass/fail; the richer diagnostics were used for failure-mode analysis and are archived as validator traces for each episode.

Inclusion/exclusion, contamination controls, and pre-execution unit-tests. The preregistered inclusion rules required that an intent pass the deterministic ambiguity detector and that mapping code successfully produce verifier inputs. We executed deterministic unit-tests for mapping code and verifier-input transformations before batch evaluation; these tests passed. Contamination checks (session-leak detection, duplicate-response entropy checks) were run; contamination\_summary.csv recorded zero flags. Deterministic reconciliation confirms episode accounting consistency and reports no provider-call failures.

Statistical analysis posture (confirmatory vs exploratory). The methodology adheres to the preregistered multiplicity and missingness rules: the primary confirmatory contrast is guarded vs baseline at \ensuremath{\alpha}=0.05 (exact McNemar two-sided) using complete\_pair\_primary (exclude incomplete pairs). Secondary contrasts were preregistered (constrained vs baseline, guarded vs constrained) with Bonferroni correction; in practice shared\_initial\_candidate semantics and adapter limits constrained the realized secondary contrasts. Analysis artifacts include paired\_contingency\_table.csv, condition\_summary.csv, and a paired bootstrap used to compute the 95\% CI (bootstrap\_resamples = 10000, bootstrap\_seed = 1729). All analysis code, seeds, and artifacts are archived for reproducibility.

Reproducibility artifacts and archival pointers. The full execution manifest, raw model-call logs, per-call runtime metadata, task-generator seeds, verifier inputs/verdicts, and analysis outputs are archived in the execution bundle referenced by the execution manifest. Deterministic reconciliation and provider accounting artifacts record model\_calls\_used = 151 and stage-level counts; these artifacts support replication of the paired accounting and the precise sampling semantics used in this trial.

\section{Results}
Execution and provider audit. The run completed as planned: completed\_episode\_count = 300 and complete\_pair\_count = 150. Provider-call accounting (audited at the provider-call level) shows model\_calls\_used = 151; completed\_hosted\_model\_call\_event\_count = 151; failed\_hosted\_model\_call\_event\_count = 0; cache\_miss\_count = 151. Stage-level counts recorded by the execution adapter were: shared\_initial\_generation = 150 and repair = 1. Contamination and mapping checks reported zero flags and no pre-execution unit-test failures.

Primary numerical outcomes and paired contingency. The paired 2\ensuremath{\times}2 contingency (guarded \ensuremath{\times} baseline) reported by the deterministic analysis executor is: n\_11 = 149 (both-pass), n\_10 = 1 (guarded-only-pass), n\_01 = 0 (baseline-only-pass), n\_00 = 0 (both-fail). Per-condition totals: baseline\_success\_count = 149/150 = 0.9933333333333333 and guarded\_success\_count = 150/150 = 1.0. The paired absolute difference (guarded - baseline) = 0.00666666666666671 (\ensuremath{\approx}0.67 percentage points). These counts are the basis for the confirmatory tests and are reproduced verbatim from the archived paired contingency artifact.

Statistical tests and bootstrap inference. The preregistered exact two-sided McNemar test on the observed discordant pair counts (n\_10 = 1, n\_01 = 0) returned p = 1.0 under the exact test implementation used by the paired\_binary\_analysis\_v1 executor. Paired bootstrap percentile 95\% confidence intervals were computed with bootstrap\_resamples = 10000 and bootstrap\_seed = 1729; the 95\% CI for the paired absolute difference is [0.00, 0.02]. These analysis parameters and results are recorded in the archived analysis artifacts and analysis log and are reported here without modification.

Repair diagnostics and revision iterations. Consistent with the near-ceiling initial success rate, only a single guarded repair invocation occurred (1 of 150 guarded episodes). That single repair converted one initially failing candidate into a passing configuration under the deterministic validator; all other guarded episodes accepted the shared initial candidate without invoking repair. Summary statistics for revision iterations in guarded episodes: median = 0, IQR = 0. Validator traces for each episode include both validation\_before and validation\_after snapshots, normalized\_configuration strings, parsed\_command\_count, required\_command\_count, observed\_touched\_field\_count, and any violation\_codes; these traces substantiate the single repair event and the distributional summary above.

Representative per-task diagnostics (artifact-grounded). Representative validator diagnostics for tasks 000001--000006 (selected from the archived representative set) illustrate variation in declared difficulty and concrete parsing outcomes while confirming consistent validator behavior:
- task-000001 (declared difficulty: medium): parsed\_command\_count = 4, required\_command\_count = 4, observed\_touched\_field\_count = 3, valid = true.
- task-000002 (declared difficulty: hard): parsed\_command\_count = 2, required\_command\_count = 2, observed\_touched\_field\_count = 2, valid = true.
- task-000003 (declared difficulty: hard): parsed\_command\_count = 6, required\_command\_count = 6, observed\_touched\_field\_count = 4, valid = true.
- task-000004 (declared difficulty: very\_hard): parsed\_command\_count = 4, required\_command\_count = 4, observed\_touched\_field\_count = 3, valid = true.
- task-000005 (declared difficulty: very\_hard): parsed\_command\_count = 3, required\_command\_count = 3, observed\_touched\_field\_count = 3, valid = true.
- task-000006 (declared difficulty: very\_hard): parsed\_command\_count = 6, required\_command\_count = 6, observed\_touched\_field\_count = 4, valid = true.
These per-task metrics and 'valid' flags are drawn directly from archived validator traces for the representative tasks and show that the verifier checked field-level counts, required sequences, and produced normalized configurations used for scoring.

Reproducibility parameters and analysis artifacts. Key analysis seeds and parameters used in inference are published in the artifacts: bootstrap\_seed = 1729 and bootstrap\_resamples = 10000. The complete paired contingency table, condition summaries, contamination summary (empty), and the full set of validator traces and raw model responses are archived. Provider accounting (model\_calls\_used = 151; repair invocations = 1) and execution manifest entries are recorded as part of the archived execution bundle to allow direct re-computation of the exact reported analyses.

Missingness, exclusions, and sensitivity checks. The preregistered missingness and exclusion rules were not triggered: failed\_hosted\_model\_call\_event\_count = 0 and unscorable episodes = 0; no ambiguity-detection exclusions, parser failures, or verifier-unsupported-feature exclusions occurred. Consequently the primary analysis used complete\_pair\_primary and did not require sensitivity treatments; the alternate preregistered sensitivity (treating missing guarded outputs as failures) was not invoked because there were no missing or failed calls. The analysis artifacts include missingness\_summary and contamination\_summary CSVs confirming these counts.

Operational interpretation of the numeric outcome. The observed baseline success rate (\ensuremath{\approx}99.33\%) leaves essentially no headroom for the preregistered detectible absolute improvement (target = +15 pp), producing a classical ceiling effect relative to the preregistered power assumptions. Under the shared\_initial\_candidate semantics used here the baseline measures the validity of the single sampled initial candidate and the guarded contrast isolates the marginal value of at most one repair applied to that identical candidate; hence, a near-perfect initial sample reduces the potential marginal utility of a single-repair guarded pipeline. The single observed repair shows the guarded pipeline can correct rare initial failures, but the aggregate effect on the preregistered estimand is numerically negligible and statistically unsupported (McNemar p = 1.0, bootstrap CI includes zero). These numeric conclusions are strictly artifact-grounded and follow directly from the archived contingency counts and analysis parameters.

Contextual diagnostics for practitioners. Beyond the binary pass/fail metric, the archived validator traces demonstrate additional operational artifacts that a guarded system would produce in deployment: normalized configuration text suitable for human audit, parsed\_command\_count and required\_command\_count to quantify task complexity, and violation\_codes that explain specific invariant violations. Even when a marginal increase in pass rate is small on a given corpus, these deterministic diagnostics are operationally useful for triage, human review, and auditable gating; the execution bundle contains representative validator traces illustrating these affordances.

\section{Discussion}
Ceiling effect and implications for preregistration. The preregistered power calculation assumed a substantially lower baseline; observed baseline\_success\_rate = 0.9933 left almost no headroom for the guarded repair to show a large absolute improvement. This ceiling effect invalidated the preregistered detectable-effect assumptions (target = 15 percentage points) and reduced the trial's ability to demonstrate benefit. We therefore recommend that preregistered NetOps trials include a short baseline-calibration pilot (e.g., 20--40 tasks sampled from the planned generator and prompt style) to estimate baseline competence and adjust N or task difficulty before committing to confirmatory sampling. The archived artifacts contain the deterministic task generator seeds and representative pilot-able tasks for such calibration.

Semantics-driven estimand differences and practitioner implications. Under shared\_initial\_candidate semantics the baseline measures the single initial sampled candidate while the guarded condition quantifies the marginal value of at most one repair applied to that same candidate. Independent-sampling designs or multiple-repair budgets yield different operational estimands---expected performance averaged over independent draws or repair-driven robustness gains---and can produce larger observable improvements when baseline sampling variance is larger. Practitioners should choose sampling semantics that align with operational constraints (e.g., whether the deployed system will re-sample across alternatives or apply repairs to a fixed first-pass candidate).

Operational affordances beyond the binary pass/fail metric. Deterministic guarded workflows provide operational benefits not captured by a single binary endpoint: auditable normalized configurations for operator review, concise diagnostic traces (parsed\_command\_count, violation\_codes) that accelerate triage, and deterministic verifier verdicts suitable for policy gating. Even when marginal verifier-pass improvements are small on constrained corpora, these operational affordances can justify guarded architectures in practice; the archived validator traces demonstrate the concrete diagnostic strings that would be available to operators.

Threats to validity. Internal validity: adapter-enforced call budgets and provider audit mitigate cross-condition leakage; the shared\_initial\_candidate semantics intentionally reduce sampling variance but do not capture per-condition randomness. Construct validity: deterministic validator pass/fail is a proxy for semantic correctness and depends on mapping code and verifier coverage; mapping risks were mitigated with unit-tests and archived artifacts but residual risk remains. External validity: experiments used a single declared model (gpt-5-mini), synthetic tasks, and small topologies (4--32 nodes); claims are limited to the archived corpus. Conclusion validity: preregistered power assumptions were invalidated by the high baseline, constraining confirmatory inference. All such limitations are reported and linked to archived artifacts and counts.

Recommendations and future experimental designs. Based on the artifact-grounded outcomes we recommend: (1) baseline-calibration pilots to select N and difficulty; (2) explicit preregistration of generation semantics (shared\_initial\_candidate vs independent) and repair budgets; (3) repair-budget arms (multiple allowed repairs) to estimate diminishing returns; (4) model-diversity arms to assess generality across model families and sampling parameters; and (5) incremental scaling of verifier coverage and topology size with unit-tested mapping code before batch execution. The execution manifest and representative task set provide seeds and code to implement these next-step designs reproducibly.

Reproducibility and artifact availability. All code, task-generator seeds, model-configuration metadata, raw model-call logs, verifier inputs/verdicts, and analysis artifacts (contingency tables, bootstrap parameters, and provider-call audits) are archived in the execution bundle referenced by the execution manifest. Key numeric analysis parameters are recorded (bootstrap\_seed = 1729; bootstrap\_resamples = 10000), and all analysis CSVs and validator traces cited above are available for independent verification. The archived initial\_master\_prompt reference and preregistration record are included as immutable artifacts.

\section{Limitations}
\begin{itemize}
\item Single-model constraint: experiments used only gpt-5-mini; results do not generalize across models or sizes.
\item Synthetic corpus and scale: tasks were programmatically generated and limited to verifier-capable toy-to-small topologies (4--32 nodes); external validity to production WANs and heterogeneous vendor CLIs is untested.
\item Shared\_initial\_candidate semantics and execution contract limited repair budget to at most one guarded repair call per intent; different generation semantics or multiple-repair budgets may yield different outcomes.
\item Ceiling effect and preregistration mismatch: baseline correctness was far higher than the pilot assumptions used for power calculations (planned detectable effect = 15 percentage points), reducing the trial's ability to demonstrate benefit and illustrating sensitivity to baseline calibration.
\item Verifier coverage and specification translation: the study measures deterministic validator pass/fail on the preregistered invariants but does not claim safety guarantees beyond the deterministic validator's modeled properties.
\item Security and adversarial robustness: the experiment did not evaluate adversarial prompt-injection or adversary-driven intents; separate targeted studies are required to assess claim-to-action vulnerabilities in agentic NetOps workflows.
\end{itemize}

\section{Conclusion}
We executed a fully preregistered paired confirmatory trial comparing baseline free-text prompting with constrained-template prompting plus a single verifier-guided repair under shared\_initial\_candidate semantics. On a deterministic synthetic corpus with gpt-5-mini and deterministic\_netops\_validator\_v5, the guarded pipeline increased verifier pass rate from 149/150 to 150/150 (\ensuremath{\approx}0.67 pp), a numerically tiny and statistically non-significant difference (exact McNemar p = 1.0; 95\% bootstrap CI = [0.00, 0.02]). The principal scientific lesson is methodological: preregistered NetOps trials should include baseline calibration to avoid ceiling effects and preserve statistical power. Technically, our artifacts bound the marginal benefit of a one-repair guarded pipeline under shared\_initial\_candidate semantics; evaluating independent-sampling semantics, larger repair budgets, model diversity, and production-scale topologies remains important future work.

\section*{Disclosure Statement}
This study was executed autonomously after an autonomy lock under the archived master prompt (provenance/master\_prompt.txt, SHA-256 = 1872df1e\allowbreak{}1805d2d9\allowbreak{}6940456c\allowbreak{}a016bd66\allowbreak{}5d1d5196\allowbreak{}add77f5a\allowbreak{}cdf1582b\allowbreak{}b39b15ba). Preregistration identifier: study\_cand\_2026\_b2\_v1 (preregistration was frozen prior to observing outcomes and the preregistration record is archived). Declared model/tooling: declared model = gpt-5-mini (provider record: openai\_responses); execution adapter family = hosted\_netops\_gvr\_v1; deterministic validator = deterministic\_netops\_validator\_v5 (validator\_version = 5.0). Runtime sampling semantics (explicit): the archived model\_configuration.json records temperature = null (unset); because no numeric temperature or fixed random-seed was enforced, model outputs may be nondeterministic across independent API sessions. Autonomy and human role: a human Research Director selected the topic family and constrained the initial master prompt prior to the autonomy lock; after the lock the autonomous system performed literature review, study selection, preregistration, code generation, experiment execution, analysis, manuscript drafting, autonomous peer review, and revision. No human manually edited technical content, analyses, figures, captions, or references after the autonomy lock. Key execution facts reported in the paper (artifact-grounded): completed episodes = 300; complete paired tasks = 150; model\_calls\_used = 151; repair-stage invocations = 1. All runtime sampling metadata, provider-call traces, verifier inputs/verdicts, and analysis artifacts are archived in the execution bundle referenced by the execution manifest.

\begin{thebibliography}{99}
\bibitem{ref1} https://doi.org/10.1002/nem.2313
\bibitem{ref2} https://doi.org/10.1145/3786583.3786898
\bibitem{ref3} https://arxiv.org/abs/2604.22513
\bibitem{ref4} https://doi.org/10.1145/3605764.3623985
\bibitem{ref5} https://doi.org/10.1145/3656454
\bibitem{ref6} https://doi.org/10.1002/widm.70111
\end{thebibliography}

\end{document}
